\section{Foundational Concepts of Model Distillation}
\label{sec:summary}

% WHAT TO WRITE -------------------------------------------------------
% Recap the problem and summarise what you have achieved so far.
% ---------------------------------------------------------------------

Model distillation serves as a method to scale down the capabilities of massive systems into formats that can be deployed on local hardware or less expensive infrastructure.

\subsection{The Teacher-Student Dynamic}

* Teacher Model: A large, high-performance LLM (e.g., DeepSeek-R1) that generates "supervision" data.
* Student Model: A smaller LLM (e.g., Qwen3 0.6B) that is fine-tuned to reproduce the teacher's reasoning path and final conclusion.

\subsection{Comparison of Distillation Types}

The documents identify three primary approaches to distillation, summarized in the table below:

Distillation Type	Training Target	Requirements	Practicality for LLMs
Hard Distillation	Teacher-generated text tokens.	Access to text outputs only.	High: Standard for LLMs; compatible with proprietary APIs.
Soft Distillation	Teacher's full probability distribution (logits).	Access to teacher logits and identical tokenizers.	Low: Computationally expensive; logits are often hidden by providers.
Combined	Both text tokens and probability distributions.	Full model access (Teacher + Student).	Medium: Common in computer vision; rarer in LLM distillation.